\chapter{Estabilidad detallada y transacciones explicadas}\label{anx:detalle}

\section{Estabilidad por escenario y arquitectura}

La Tabla \ref{tab:esc_arq} desglosa por arquitectura la estabilidad por escenario de la Tabla
\ref{tab:h1}, con pesos por clase. Cada celda es la media de tres semillas. Incluye GCN, que no
supera el filtro de calidad con esta pérdida. Ninguna arquitectura muestra una caída sistemática
hacia 1:100 y 1:200.

\begin{table}[ht]
\centering
\caption{Estabilidad entre réplicas por escenario y arquitectura (pesos por clase, media de 3
semillas).}
\label{tab:esc_arq}
\scriptsize
\setlength{\tabcolsep}{3pt}
% Generado por tesis_v4/scripts/analisis.py. No editar a mano.
\begin{tabular}{lllcccccc}
\toprule
Medida & Explicador & Arquitectura & 1:10 & 1:10 + SMOTE & 1:20 & Nativo & 1:100 & 1:200 \\
\midrule
Spearman global & GNNExplainer & GCN & 0{,}936 & 0{,}937 & 0{,}937 & 0{,}937 & 0{,}936 & 0{,}942 \\
 &  & GraphSAGE & 0{,}929 & 0{,}926 & 0{,}953 & 0{,}916 & 0{,}957 & 0{,}937 \\
 &  & GAT & 0{,}937 & 0{,}919 & 0{,}918 & 0{,}918 & 0{,}918 & 0{,}921 \\
 &  & TAGCN & 0{,}940 & 0{,}943 & 0{,}917 & 0{,}917 & 0{,}943 & 0{,}938 \\
 & Shapley & GCN & 0{,}996 & 0{,}996 & 0{,}996 & 0{,}996 & 0{,}996 & 0{,}996 \\
 &  & GraphSAGE & 0{,}994 & 0{,}994 & 0{,}992 & 0{,}995 & 0{,}992 & 0{,}992 \\
 &  & GAT & 0{,}994 & 0{,}993 & 0{,}993 & 0{,}993 & 0{,}993 & 0{,}992 \\
 &  & TAGCN & 0{,}989 & 0{,}988 & 0{,}993 & 0{,}992 & 0{,}988 & 0{,}987 \\
 & Expected Gradients & GCN & 0{,}923 & 0{,}922 & 0{,}923 & 0{,}923 & 0{,}918 & 0{,}914 \\
 &  & GraphSAGE & 0{,}941 & 0{,}946 & 0{,}906 & 0{,}948 & 0{,}915 & 0{,}945 \\
 &  & GAT & 0{,}939 & 0{,}930 & 0{,}929 & 0{,}929 & 0{,}929 & 0{,}924 \\
 &  & TAGCN & 0{,}912 & 0{,}915 & 0{,}926 & 0{,}926 & 0{,}911 & 0{,}913 \\
\midrule
Spearman $\neq 0$ & GNNExplainer & GCN & 0{,}689 & 0{,}695 & 0{,}690 & 0{,}690 & 0{,}688 & 0{,}720 \\
 &  & GraphSAGE & 0{,}660 & 0{,}651 & 0{,}781 & 0{,}596 & 0{,}800 & 0{,}701 \\
 &  & GAT & 0{,}692 & 0{,}616 & 0{,}616 & 0{,}615 & 0{,}615 & 0{,}626 \\
 &  & TAGCN & 0{,}688 & 0{,}701 & 0{,}565 & 0{,}566 & 0{,}706 & 0{,}674 \\
 & Shapley & GCN & 0{,}993 & 0{,}994 & 0{,}994 & 0{,}993 & 0{,}993 & 0{,}993 \\
 &  & GraphSAGE & 0{,}992 & 0{,}991 & 0{,}987 & 0{,}993 & 0{,}989 & 0{,}989 \\
 &  & GAT & 0{,}992 & 0{,}990 & 0{,}989 & 0{,}989 & 0{,}989 & 0{,}987 \\
 &  & TAGCN & 0{,}980 & 0{,}980 & 0{,}990 & 0{,}989 & 0{,}978 & 0{,}978 \\
 & Expected Gradients & GCN & 0{,}958 & 0{,}958 & 0{,}958 & 0{,}957 & 0{,}955 & 0{,}952 \\
 &  & GraphSAGE & 0{,}968 & 0{,}973 & 0{,}943 & 0{,}975 & 0{,}951 & 0{,}960 \\
 &  & GAT & 0{,}965 & 0{,}959 & 0{,}959 & 0{,}958 & 0{,}959 & 0{,}957 \\
 &  & TAGCN & 0{,}944 & 0{,}943 & 0{,}961 & 0{,}961 & 0{,}946 & 0{,}945 \\
\midrule
Jaccard top-10 & GNNExplainer & GCN & 0{,}703 & 0{,}709 & 0{,}709 & 0{,}710 & 0{,}713 & 0{,}699 \\
 &  & GraphSAGE & 0{,}604 & 0{,}605 & 0{,}669 & 0{,}611 & 0{,}660 & 0{,}600 \\
 &  & GAT & 0{,}656 & 0{,}645 & 0{,}648 & 0{,}650 & 0{,}647 & 0{,}648 \\
 &  & TAGCN & 0{,}602 & 0{,}617 & 0{,}526 & 0{,}526 & 0{,}605 & 0{,}552 \\
 & Shapley & GCN & 0{,}946 & 0{,}949 & 0{,}948 & 0{,}950 & 0{,}940 & 0{,}950 \\
 &  & GraphSAGE & 0{,}936 & 0{,}921 & 0{,}913 & 0{,}944 & 0{,}901 & 0{,}900 \\
 &  & GAT & 0{,}943 & 0{,}915 & 0{,}924 & 0{,}924 & 0{,}926 & 0{,}911 \\
 &  & TAGCN & 0{,}873 & 0{,}860 & 0{,}925 & 0{,}925 & 0{,}843 & 0{,}847 \\
 & Expected Gradients & GCN & 0{,}797 & 0{,}803 & 0{,}800 & 0{,}792 & 0{,}796 & 0{,}794 \\
 &  & GraphSAGE & 0{,}808 & 0{,}809 & 0{,}776 & 0{,}805 & 0{,}777 & 0{,}830 \\
 &  & GAT & 0{,}828 & 0{,}818 & 0{,}807 & 0{,}812 & 0{,}812 & 0{,}813 \\
 &  & TAGCN & 0{,}813 & 0{,}815 & 0{,}800 & 0{,}794 & 0{,}791 & 0{,}824 \\
\bottomrule
\end{tabular}

\end{table}

\section{Las transacciones explicadas}

La Tabla \ref{tab:nodos} describe las 30 transacciones ilícitas de validación sobre las que se
midió la estabilidad. Todas pertenecen a los pasos de tiempo 30 a 35, tienen pocas aristas y
entre la cuarta parte y la mitad de sus variables valen cero. En promedio, cada una es
clasificada correctamente por el 64,5~\% de los 54 modelos con pesos por clase de las tres
arquitecturas que aprenden, de modo que el conjunto incluye transacciones que casi todos los
modelos detectan y otras que casi ninguno detecta.

{\small
% Generado por tesis_v4/scripts/analisis.py. No editar a mano.
\begin{longtable}{rrrrc}
\caption{Las 30 transacciones ilícitas de validación explicadas. Nodo: índice en el grafo; Paso: paso de tiempo; Grado: número de aristas; Nulas: fracción de variables que valen 0 tras el escalado; Acierto: fracción de los 54 modelos con pesos por clase (GraphSAGE, GAT y TAGCN; 6 escenarios; 3 semillas) que la clasifican como ilícita con umbral 0,5.}\label{tab:nodos}\\
\toprule
Nodo & Paso & Grado & Nulas & Acierto \\
\midrule
\endfirsthead
\toprule
Nodo & Paso & Grado & Nulas & Acierto \\
\midrule
\endhead
123122 & 30 & 1 & 0{,}50 & 0{,}69 \\
123491 & 31 & 2 & 0{,}23 & 0{,}41 \\
123935 & 31 & 2 & 0{,}48 & 1{,}00 \\
124436 & 31 & 2 & 0{,}42 & 0{,}26 \\
125970 & 31 & 1 & 0{,}46 & 0{,}19 \\
126209 & 32 & 2 & 0{,}31 & 0{,}56 \\
126246 & 32 & 2 & 0{,}50 & 0{,}70 \\
126320 & 32 & 2 & 0{,}39 & 0{,}52 \\
126853 & 32 & 2 & 0{,}39 & 0{,}50 \\
126854 & 32 & 2 & 0{,}42 & 0{,}52 \\
127435 & 32 & 1 & 0{,}39 & 0{,}54 \\
128100 & 32 & 2 & 0{,}42 & 0{,}41 \\
128946 & 32 & 1 & 0{,}48 & 0{,}43 \\
128985 & 32 & 1 & 0{,}47 & 0{,}56 \\
129303 & 32 & 2 & 0{,}45 & 0{,}89 \\
129889 & 32 & 2 & 0{,}42 & 0{,}52 \\
130250 & 32 & 1 & 0{,}45 & 0{,}98 \\
130493 & 32 & 2 & 0{,}48 & 1{,}00 \\
136391 & 35 & 2 & 0{,}47 & 0{,}74 \\
136551 & 35 & 1 & 0{,}47 & 0{,}65 \\
136736 & 35 & 2 & 0{,}46 & 0{,}65 \\
139218 & 35 & 2 & 0{,}42 & 1{,}00 \\
139219 & 35 & 2 & 0{,}48 & 0{,}72 \\
139715 & 35 & 1 & 0{,}50 & 0{,}31 \\
139737 & 35 & 1 & 0{,}50 & 0{,}89 \\
140788 & 35 & 1 & 0{,}42 & 0{,}44 \\
140789 & 35 & 1 & 0{,}47 & 0{,}98 \\
140830 & 35 & 1 & 0{,}47 & 0{,}65 \\
140890 & 35 & 1 & 0{,}48 & 0{,}98 \\
140978 & 35 & 1 & 0{,}46 & 0{,}69 \\
\bottomrule
\end{longtable}

}
